\documentclass[11pt]{article}
\usepackage[letterpaper,margin=1in]{geometry}
\usepackage[T1]{fontenc}
\usepackage{lmodern,amsmath,amssymb,amsthm,mathtools,microtype}
\usepackage[hidelinks]{hyperref}
\hypersetup{pdftitle={An explicit weighted localization bound for a signed prime-gamma form},pdfauthor={Jongmin Choi}}
\newtheorem{theorem}{Theorem}[section]
\newtheorem{proposition}[theorem]{Proposition}
\newtheorem{lemma}[theorem]{Lemma}
\newtheorem{corollary}[theorem]{Corollary}
\theoremstyle{remark}\newtheorem{remark}[theorem]{Remark}
\newcommand{\R}{\mathbb R}
\newcommand{\dd}{\,\mathrm d}
\newcommand{\norm}[1]{\lVert#1\rVert}
\newcommand{\Pole}{\operatorname{Pole}}
\newcommand{\supp}{\operatorname{supp}}
\newcommand{\HH}{\mathcal H}
\numberwithin{equation}{section}
\allowdisplaybreaks[2]
\setlength{\emergencystretch}{2em}
\title{An explicit weighted localization bound\\for a signed prime--gamma form}
\author{Jongmin Choi\\[3pt]\small Independent Researcher, Seoul, Korea\\\small ORCID \href{https://orcid.org/0009-0008-7448-514X}{0009-0008-7448-514X}}
\date{27 October 2026\\\small Preprint version 1.0}
\begin{document}
\maketitle
\begin{abstract}
We study a real even quadratic form consisting of a nonlocal gamma energy,
the complete prime-power correlation sum, a local mass term, and a rank-one
pole term. In an explicit Jacobi-theta weighted space, we construct a cubic
angle partition of unity and bound the complete signed localization defect
without losing any exterior gamma energy. For an interior radius $R\ge1$
and $S=e^{2R}$, the defect is at most
$9000(1+\log S)S^{-5/2}e^{-\pi S/2}$ times the squared weighted norm.
We also cancel the deterministic finite-prime main phase against the full
pole exactly and prove a complete weighted remainder estimate. Matching
the cutoff support to $\log N$ gives a finite-$N$ localization formula with
no second interior truncation error. The theta inputs and all constants are
proved directly. The remaining interior form is infinite-dimensional and
signed; no positivity theorem or conclusion about the Riemann hypothesis
is obtained.
\end{abstract}
\begingroup\small
\noindent\textbf{Keywords.} Weighted quadratic forms; nonlocal localization;
prime powers; Jacobi theta function; signed correlation.\par
\noindent\textbf{License.} Copyright \textcopyright\ 2026 Jongmin Choi.
This work is licensed under Creative Commons Attribution--NonCommercial
4.0 International (CC BY-NC 4.0).\par
\endgroup

\section{The form and the localization statement}

The issue addressed here is quantitative localization, not the sign of the
localized form. An exact partition-of-unity expansion leaves a signed
cross-edge defect. The rapid decay of a theta weight makes that defect
small, and a cubic angle taper supplies an additional fourth-order gain.
The proof keeps the complete prime-power sum and the pole term throughout.
Generic partition identities and exterior localization are the underlying
mechanisms; no historical novelty or priority claim is made for them.

For $t>0$ define
\[
 \vartheta(t)=\sum_{j\in\mathbb Z}e^{-\pi j^2t},\qquad
 G(x)=e^{x/2}\vartheta(e^{2x}),\qquad
 \kappa(x)=\tfrac12\bigl(G''(x)-G(x)/4\bigr).
\]
Put $C(x)=\cosh(x/2)$, $W(x)=\kappa(x)/(2C(x))$, and
\begin{equation}\label{eq:X}
 X=L^2_{\mathrm{even}}(\R,W^{-1}\dd x;\R),\qquad
 N_X(u)=\norm{u}_X^2=\int_\R\frac{|u(x)|^2}{W(x)}\dd x.
\end{equation}
Appendix~\ref{app:theta} proves, independently of any zero-location
statement, that $W$ is smooth, positive and even, and that
\begin{equation}\label{eq:theta-inputs}
 W(x)\le40\exp(4|x|-\pi e^{2|x|}),\qquad
 \int_\R2C(x)\kappa(x)\dd x=1.
\end{equation}

Initially let $u$ belong to the real even core
$\mathcal D_0=C_c^\infty(\R;\R)\cap X$. We use
\begin{equation}\label{eq:fourier-correlation}
 F_u(\xi)=\int_\R u(x)e^{-i\xi x}\dd x,\qquad
 H_u(s)=\int_\R u(x)u(x+s)\dd x,
 \qquad M(u)=\norm{u}_2^2.
\end{equation}
The correlation is real and even. With $\Lambda$ denoting the von Mangoldt
function, so that $\Lambda(p^k)=\log p$ for every prime $p$ and integer
$k\ge1$ and $\Lambda(n)=0$ otherwise, set
\begin{align}
 \rho(s)&=\frac{e^{-s/2}}{1-e^{-2s}},&
 t_\Gamma(u)&=\int_0^\infty\rho(s)\norm{u(\cdot+s)-u}_2^2\dd s,
 \label{eq:gamma}\\
 a_n&=\frac{\Lambda(n)}{\sqrt n},&
 p_N(u)&=2\sum_{2\le n\le N}a_nH_u(\log n),
 \label{eq:prime}\\
 c_\Gamma&=\log\pi+\gamma+\frac\pi2+3\log2,&
 \Pole(u)&=2\left|\int_\R C(x)u(x)\dd x\right|^2.
 \label{eq:constant-pole}
\end{align}
Here $\gamma$ is Euler's constant. The elementary bounds
$0<\gamma<1$, $\log\pi<2$, $\pi/2<2$ and $\log2<1$ give
$0<c_\Gamma<8$. The first bounds follow by comparing harmonic sums
with the integral of $1/t$; the others follow from $3<\pi<4$ and
$2<e<3$. The pole equals $2|F_u(i/2)|^2$ because $u$ is real and even.
Define
\begin{equation}\label{eq:forms}
 q_N(u)=t_\Gamma(u)-c_\Gamma M(u)-p_N(u)+\Pole(u),
 \qquad q_\infty(u)=t_\Gamma(u)-c_\Gamma M(u)-p_\infty(u)+\Pole(u).
\end{equation}
The complete series $p_\infty$ is an $X$-bounded form, as proved below.
The gamma energy is once-oriented: its shifts run over $s>0$ only.
The factor two in \eqref{eq:prime} retains both prime orientations.

Our cutoffs are completely specified. Let
\begin{equation}\label{eq:cutoffs}
 b(t)=\begin{cases}0,&t\le0,\\3t^2-2t^3,&0<t<1,\\1,&t\ge1,\end{cases}
 \quad \theta_R(x)=\frac\pi2b(|x|-R),\quad
 \chi_R(x)=\cos\theta_R(x),\quad\eta_R(x)=\sin\theta_R(x).
\end{equation}
They are even $C^1$ functions, not necessarily smooth at their transition
endpoints. They satisfy $\chi_R^2+\eta_R^2=1$, $\chi_R=1$ on $[-R,R]$
and $\chi_R=0$ outside $[-R-1,R+1]$. Section~\ref{sec:domain} admits their
products with core tests by compact $H^1$ approximation.

\begin{theorem}[Explicit signed localization]\label{thm:main}
For every $R\ge1$ and $u\in\mathcal D_0$,
\begin{equation}\label{eq:main}
 \left|q_\infty(u)-q_\infty(\chi_Ru)-t_\Gamma(\eta_Ru)\right|
 \le\varepsilon_RN_X(u),
\end{equation}
where, with $S=e^{2R}$,
\begin{equation}\label{eq:closed-eps}
 \varepsilon_R\le9000(1+\log S)S^{-5/2}e^{-\pi S/2}.
\end{equation}
An explicit, sharper expression for $\varepsilon_R$ is given in
\eqref{eq:eps}. In particular the coefficient of the retained exterior
energy $t_\Gamma(\eta_Ru)$ is exactly one.
\end{theorem}

The conclusion is a vanishing error in the original weighted norm. The
single uncontrolled term $q_\infty(\chi_Ru)$ has finite physical support
but is still infinite-dimensional. A finite number of prime shifts on
that support does not reduce its gamma energy to a finite-dimensional
problem. Theorem~\ref{thm:main} does not assert that this term is positive.

\section{Weighted correlation bounds and the form domain}\label{sec:domain}

Write $\beta=\sqrt W$. The envelope in \eqref{eq:theta-inputs} implies
that $X$ embeds continuously in ordinary $L^2$ and that every $u\in X$
has every finite exponential $L^1$ moment: for $a\ge0$,
\[
 \int_\R e^{a|x|}|u(x)|\dd x
 \le N_X(u)^{1/2}\left(\int_\R e^{2a|x|}W(x)\dd x\right)^{1/2}<\infty.
\]
Consequently all pole and correlation integrals used below are absolutely
justified on $X$.

\begin{lemma}[Complete weighted tails]\label{lem:tails}
For $s\ge0$ and $u\in X$,
\begin{equation}\label{eq:overlap}
 \beta(x)\beta(x+s)\le h(s):=40e^{2s-\pi e^s},\qquad
 |H_u(s)|\le h(s)N_X(u).
\end{equation}
For every integer $N\ge1$,
\begin{equation}\label{eq:prime-tail}
 |p_\infty(u)-p_N(u)|\le r_NN_X(u),\qquad
 r_N=80\sum_{n>N}(\log n)n^{3/2}e^{-\pi n}.
\end{equation}
Moreover $|p_\infty(u)|\le r_1N_X(u)$ and
\begin{equation}\label{eq:pole-bound}
 0\le\Pole(u)\le\tfrac12N_X(u).
\end{equation}
For $A\ge2$ define
\begin{equation}\label{eq:tau}
 \tau_A=80\int_A^\infty t^{3/2}e^{-\pi t}\dd t
 \le\frac{80A^{3/2}e^{-\pi A}}{\pi-3/(2A)}.
\end{equation}
\end{lemma}
\begin{proof}
Put $a=|x|$, $b=|x+s|$ and $v=a+b\ge s$. The arithmetic--geometric
mean inequality and the decrease of $2v-\pi e^v$ for $v\ge0$ give
\[
 \beta(x)\beta(x+s)
 \le40e^{2v-(\pi/2)(e^{2a}+e^{2b})}
 \le40e^{2v-\pi e^v}\le40e^{2s-\pi e^s}.
\]
Writing $u=\beta f$ and applying ordinary translation-invariant
Cauchy--Schwarz proves the correlation estimate. At $s=\log n$, use
$\Lambda(n)\le\log n$ and sum every $n>N$ to obtain
\eqref{eq:prime-tail}. No prime power is omitted. Finally
\[
 \int_\R C^2W\dd x=\frac12\int_\R C\kappa\dd x=\frac14,
\]
which proves \eqref{eq:pole-bound}. Concavity of $\log t$ gives
$t^{3/2}\le A^{3/2}e^{3(t-A)/(2A)}$ for $t\ge A$; integration proves
\eqref{eq:tau}.
\end{proof}

The Fourier multiplier of the gamma form is
\begin{equation}\label{eq:d-gamma}
 d_\Gamma(\xi)=2\int_0^\infty\rho(s)(1-\cos(\xi s))\dd s,
 \qquad t_\Gamma(u)=\frac1{2\pi}\int_\R d_\Gamma(\xi)|F_u(\xi)|^2\dd\xi.
\end{equation}
Tonelli and Plancherel justify the nonnegative integral. Expanding
$\rho(s)=\sum_{j\ge0}e^{-(2j+1/2)s}$ gives
\[
 0\le d_\Gamma(\xi)\le\xi^2\int_0^\infty s^2\rho(s)\dd s
 =2\xi^2\sum_{j\ge0}(2j+1/2)^{-3}\le18\xi^2.
\]
Indeed the $j=0$ term is $16$, while the others sum to at most
$\frac14\sum_{j\ge1}j^{-3}\le3/8$.

If $v$ is real, even, compactly supported and in $H^1(\R)$, convolution
with even compact smooth approximate identities gives $v_j\in\mathcal D_0$
converging in $H^1$, with supports in one common compact interval. Since
$W$ has a positive minimum on that interval, $v_j\to v$ in $X$, and
\begin{equation}\label{eq:admission}
 t_\Gamma(v_j-v)\le18\norm{(v_j-v)'}_2^2\longrightarrow0.
\end{equation}
The prime, pole and local mass terms are continuous in $X$. Thus all form
values converge. This admits $\chi_Ru$ and $\eta_Ru$ to the completion of
the original core in the norm $(N_X+t_\Gamma)^{1/2}$. No identification
with an unspecified maximal form domain is required. In particular the
gamma form itself has not been declared bounded on $X$.

\section{Exact main-phase cancellation}\label{sec:phase}

For an integer $N\ge2$, put $L=\log N$ and
\begin{align}
 P_N(\xi)&=2\sum_{2\le n\le N}a_n\cos(\xi\log n),\label{eq:PN}\\
 M_N(\xi)&=2\operatorname{Re}\frac{N^{1/2+i\xi}}{1/2+i\xi}
 =\frac{4\sqrt N\cos(L\xi-\arctan(2\xi))}{\sqrt{1+4\xi^2}},
 \qquad\Delta_N=P_N-M_N.\label{eq:MN}
\end{align}
Both finite-$N$ multipliers are bounded and real even. Define
\begin{equation}\label{eq:renormalized}
 \HH_N(u)=\frac1{2\pi}\int_\R
       (d_\Gamma-c_\Gamma-\Delta_N)|F_u|^2\dd\xi
\end{equation}
on the compact core and its admitted compact $H^1$ products.

\begin{proposition}[Full pole and complete remainder]\label{prop:phase}
For $u\in X$,
\begin{equation}\label{eq:pole-main}
 \Pole(u)-\frac1{2\pi}\int_\R M_N|F_u|^2\dd\xi=T_N(u),\qquad
 T_N(u)=2\int_{\log N}^\infty e^{s/2}H_u(s)\dd s,
\end{equation}
and $|T_N(u)|\le\tau_NN_X(u)$. On the compact core,
\begin{equation}\label{eq:phase-form}
 q_N(u)=\HH_N(u)+T_N(u),\qquad
 |\HH_N(u)-q_\infty(u)|\le(r_N+\tau_N)N_X(u).
\end{equation}
The discrepancy forms satisfy, on all of $X$,
\begin{equation}\label{eq:discrepancy}
 \left|\frac1{2\pi}\int_\R\Delta_N|F_u|^2\dd\xi
          -\bigl(p_\infty(u)-\Pole(u)\bigr)\right|
 \le(r_N+\tau_N)N_X(u).
\end{equation}
\end{proposition}
\begin{proof}
Absolute Fubini and $y=x+s$ give
\begin{equation}\label{eq:full-pole}
 \int_\R e^{s/2}H_u(s)\dd s
 =\left(\int_\R e^{-x/2}u(x)\dd x\right)
  \left(\int_\R e^{y/2}u(y)\dd y\right)=|F_u(i/2)|^2.
\end{equation}
The lower endpoint is $-\infty$, not zero. The negative half-line supplies
a genuine term and cannot be discarded. On the other hand,
\[
 M_N(\xi)=2\operatorname{Re}\int_{-\infty}^L e^{(1/2+i\xi)s}\dd s.
\]
The last kernel is integrable for fixed $N$, and $|F_u|^2\in L^1$.
Fourier inversion for the correlation and Fubini therefore yield
\[
 \frac1{2\pi}\int_\R M_N|F_u|^2\dd\xi
       =2\int_{-\infty}^L e^{s/2}H_u(s)\dd s.
\]
Subtracting from twice \eqref{eq:full-pole} proves
\eqref{eq:pole-main}. The overlap bound and $t=e^s$ give
\[
 |T_N(u)|\le80\int_L^\infty e^{5s/2-\pi e^s}\dd s\,N_X(u)
 =\tau_NN_X(u).
\]
Substitution into \eqref{eq:forms} proves the first identity in
\eqref{eq:phase-form}. Also
\[
 \frac1{2\pi}\int_\R\Delta_N|F_u|^2\dd\xi
       =p_N(u)-\Pole(u)+T_N(u),
\]
so \eqref{eq:prime-tail} proves the remaining statements.
\end{proof}

Polarization and the Riesz representation theorem express
\eqref{eq:discrepancy} as bounded selfadjoint operator convergence on the
real Hilbert space $X$, with norm error at most $r_N+\tau_N$. There is no
extra polarization factor: for a bounded symmetric operator the norm is
the supremum of the absolute diagonal pairing over unit vectors. This
does not bound $\int|\Delta_N||F_u|^2$, and the limiting signed operator
$p_\infty-\Pole$ is not claimed to have vanishing norm. In
\eqref{eq:phase-form} only the difference of the two gamma-containing
forms is $X$-bounded. The full pole has been consumed once by
\eqref{eq:pole-main}; it is not an additional resource after that identity.

For clarity, the arithmetic endpoint convention can also be read directly.
Let $\Psi(t)=\sum_{n\le t}\Lambda(n)$ be right-continuous and
$E(t)=\Psi(t)-t$, so $E(1)=-1$. Stieltjes integration by parts, including
an atom at $N$, gives for $z=1/2+i\xi$
\begin{equation}\label{eq:abel}
 \sum_{n\le N}\Lambda(n)n^{z-1}
 =\frac{N^z}{z}+\frac{z-1}{z}+E(N)N^{z-1}
       -(z-1)\int_1^N E(t)t^{z-2}\dd t.
\end{equation}
Thus $\Delta_N$ is twice the real part of the last three terms. The lower
endpoint $(z-1)/z$ is present. No estimate for $E$ or prime number theorem
is used in Proposition~\ref{prop:phase} or in the localization theorem.

\section{The weighted cross-edge gain}\label{sec:cross-edge}

The cubic in \eqref{eq:cutoffs} obeys $b'(t)=6t(1-t)$ on $(0,1)$.
Consequently
\begin{equation}\label{eq:angle-bounds}
 0\le\theta_R\le\pi/2,\qquad
 \operatorname{Lip}(\theta_R)\le3\pi/4,\qquad
 \theta_R(x)\le\frac{3\pi}{2}(|x|-R)_+^2.
\end{equation}
For $x,y\in\R$ define
\begin{align}
 \ell_R(x,y)&=1-\chi_R(x)\chi_R(y)-\eta_R(x)\eta_R(y)
             =1-\cos(\theta_R(x)-\theta_R(y)),\label{eq:ell}\\
 e_R(x,y)&=1-\chi_R(x)\chi_R(y).\label{eq:edge-e}
\end{align}
Both lie in $[0,1]$ and vanish if both endpoints lie in $[-R,R]$.
Use the following constants throughout the proof:
\begin{equation}\label{eq:constants}
 \begin{gathered}
 S=e^{2R},\qquad k_R=\pi S-2,\qquad
 G_R=40S e^{-\pi(S+1)/2},\\
 B_R=\frac{288\pi^2}{e^4k_R^4}G_R,\qquad
 \alpha_R=40S^2e^{-\pi S}.
 \end{gathered}
\end{equation}

\begin{lemma}[Fourth-order taper bound]\label{lem:edges}
For all $x,y\in\R$ and $R\ge1$,
\begin{equation}\label{eq:edge-bounds}
 \ell_R(x,y)\beta(x)\beta(y)\le B_R,\qquad
 e_R(x,y)\beta(x)\beta(y)\le2B_R.
\end{equation}
For $s\ge0$ one also has
\begin{equation}\label{eq:near-diagonal}
 \ell_R(x,x+s)\beta(x)\beta(x+s)
 \le\min\left\{\frac{9\pi^2}{32}G_Rs^2,B_R\right\}.
\end{equation}
Either defect times $\beta(x)\beta(x+s)$ is at most $h(s)$.
\end{lemma}
\begin{proof}
If a defect is nonzero, set $\max(|x|,|y|)=R+\delta$ with $\delta>0$.
The function
\[
 a(t)=\sqrt{40}\exp\bigl(2t-(\pi/2)e^{2t}\bigr),\qquad t\ge0,
\]
is decreasing and has concave logarithm. Since $\beta(x)\le a(|x|)$,
the tangent bound at $R$ gives
\begin{equation}\label{eq:edge-envelope}
 \beta(x)\beta(y)\le a(R+\delta)a(0)
                    \le G_Re^{-k_R\delta}.
\end{equation}
Monotonicity of the angle in $|x|$ and $1-\cos t\le t^2/2$ give
\[
 \ell_R(x,y)\le\frac{9\pi^2}{8}\delta^4,\qquad
 e_R(x,y)\le\frac{\theta_R(x)^2+\theta_R(y)^2}{2}
                  \le\frac{9\pi^2}{4}\delta^4.
\]
Now $\sup_{\delta\ge0}\delta^4e^{-k_R\delta}=256/(e^4k_R^4)$,
which proves \eqref{eq:edge-bounds}. The Lipschitz estimate in
\eqref{eq:angle-bounds} gives
$\ell_R(x,x+s)\le9\pi^2s^2/32$; on a nonzero edge
\eqref{eq:edge-envelope} also bounds the product by $G_R$.
This proves \eqref{eq:near-diagonal}. The final assertion follows from
\eqref{eq:overlap} and the bound of one on either defect.
\end{proof}

The factor $k_R^{-4}$ records the order of the taper at the inner endpoint.
The exponential factor in $G_R$ comes from the actual large endpoint of
the edge. Neither argument uses a sampled source value.

\section{Exact signed localization and complete error accounting}

For compact $H^1$ inputs define
\[
 I_\ell(s)=\int_\R\ell_R(x,x+s)u(x)u(x+s)\dd x,\qquad
 I_e(s)=\int_\R e_R(x,x+s)u(x)u(x+s)\dd x.
\]
Expansion of the gamma squares gives exactly
\begin{equation}\label{eq:gamma-localization}
 t_\Gamma(\chi_Ru)+t_\Gamma(\eta_Ru)-t_\Gamma(u)
                 =2\int_0^\infty\rho(s)I_\ell(s)\dd s.
\end{equation}
Evenness in \eqref{eq:full-pole} yields
\begin{equation}\label{eq:positive-half-pole}
 \Pole(u)=4\int_0^\infty C(s)H_u(s)\dd s.
\end{equation}
The same identity holds separately for the cutoff products because the
cutoffs are even. Subtracting the full forms therefore gives the direct
one-cutoff identity
\begin{equation}\label{eq:exact-localization}
 \begin{aligned}
 q_\infty(u)-q_\infty(\chi_Ru)-t_\Gamma(\eta_Ru)
 ={}&-c_\Gamma\norm{\eta_Ru}_2^2
       -2\int_0^\infty\rho(s)I_\ell(s)\dd s\\
    &-2\sum_{n\ge2}a_nI_e(\log n)
       +4\int_0^\infty C(s)I_e(s)\dd s.
 \end{aligned}
\end{equation}
In particular the last term is the signed difference
$\Pole(u)-\Pole(\chi_Ru)$, not an extra positive allowance.

Write $u=\beta f$, so $\norm{f}_2^2=N_X(u)$. Each uniform weighted
edge bound becomes a correlation bound with the same constant by
$\int|f(x)f(x+s)|\dd x\le\norm{f}_2^2$. We estimate the four terms
in \eqref{eq:exact-localization} separately, without splitting any
frequency observation into positive and negative parts.

\subsection{The gamma defect}
Put
\begin{equation}\label{eq:gamma-split}
 s_R^*=\frac{32}{e^2k_R^2}<1,\qquad
 L_R^*=\frac52+\log\frac{e^2k_R^2}{32}.
\end{equation}
For $0<s\le1$, $\rho(s)=e^{s/2}/(2\sinh s)\le1/s$ because
$e^{1/2}<2$ and $\sinh s\ge s$. Also
\[
 \int_1^\infty\rho(s)\dd s
 \le\frac{2e^{-1/2}}{1-e^{-2}}<2;
\]
the last inequality follows from $e^{1/2}>3/2$ and $e^2>4$.
The two alternatives in \eqref{eq:near-diagonal} are equal at $s_R^*$.
Splitting at $s_R^*$ and at $1$ gives
\begin{equation}\label{eq:gamma-error}
 2\left|\int_0^\infty\rho(s)I_\ell(s)\dd s\right|
 \le2B_R\left(\frac12+\log\frac1{s_R^*}+2\right)N_X(u)
 =2B_RL_R^*N_X(u).
\end{equation}

\subsection{Every prime power and the full pole}
Let $K=\lfloor S\rfloor$. Since $\Lambda(n)\le\log n$,
\[
 \sum_{2\le n\le S}a_n
 \le(\log S)\sum_{2\le n\le S}n^{-1/2}
 \le2(\log S)\sqrt S=4Re^R.
\]
Use $2B_R$ for the interior edge bound and $h(\log n)$ for all $n>S$.
Then
\begin{equation}\label{eq:prime-error}
 2\left|\sum_{n\ge2}a_nI_e(\log n)\right|
 \le\bigl(16Re^RB_R+r_K\bigr)N_X(u).
\end{equation}
The strict tail is exactly $n>K$, including every prime power there.

For the pole defect, use $2B_R$ on $0\le s\le2R$ and the full overlap
envelope on the remaining ray. Since $\int_0^{2R}C(s)\dd s=2\sinh R$,
the first contribution is at most $8e^RB_RN_X(u)$. In the tail,
$4C(s)h(s)=80(e^{5s/2}+e^{3s/2})e^{-\pi e^s}$.
With $t=e^s$, the growing part gives $\tau_S$, while the other is at
most $\tau_S/S$. Therefore
\begin{equation}\label{eq:pole-error}
 4\left|\int_0^\infty C(s)I_e(s)\dd s\right|
 \le\bigl[8e^RB_R+(1+S^{-1})\tau_S\bigr]N_X(u).
\end{equation}
Finally $\eta_R=0$ on $[-R,R]$ and the envelope decreases for $|x|\ge0$,
so $c_\Gamma\norm{\eta_Ru}_2^2\le8\alpha_RN_X(u)$.

Combining these estimates proves \eqref{eq:main} with
\begin{equation}\label{eq:eps}
 \boxed{\displaystyle
 \varepsilon_R=2B_R\bigl[L_R^*+(8R+4)e^R\bigr]
       +8\alpha_R+r_{\lfloor S\rfloor}+(1+S^{-1})\tau_S.}
\end{equation}
This estimates the direct identity \eqref{eq:exact-localization}. No
standalone exterior error is added a second time. In particular no fraction
of $t_\Gamma(u)$ has been spent in the error bound.

\subsection{An elementary bound for all residual constants}
For $n\ge1$, $\log n\le\sqrt n$: the maximum of
$\log t/\sqrt t$ on $t\ge1$ is $2/e<1$. Let $q=e^{-\pi}<1/8$.
The geometric second moment, with $m=K+1$, yields
\[
 \sum_{n>K}n^2q^n
 =m^2q^m\left\{\frac1{1-q}
       +\frac{2q}{m(1-q)^2}
       +\frac{q(1+q)}{m^2(1-q)^3}\right\}
 <2m^2q^m,
\]
since the bracket is at most
$8/7+16/49+72/343=576/343<2$. Hence
\begin{equation}\label{eq:r-closed}
 r_K\le160(K+1)^2e^{-\pi(K+1)}\qquad(K\ge1).
\end{equation}
Because $S\ge e^2>4$, \eqref{eq:tau}, \eqref{eq:constants} and
\eqref{eq:r-closed} imply
\begin{equation}\label{eq:residual-closed}
 8\alpha_R+r_{\lfloor S\rfloor}+(1+S^{-1})\tau_S
 <(320+250+25)S^2e^{-\pi S}<600S^2e^{-\pi S}.
\end{equation}
For the middle term use $K+1\le S+1<5S/4$ and $K+1>S$.
For the last use $\pi-3/(2S)>21/8$, $1+1/S<5/4$ and $\sqrt S>2$.

Also $k_R>(5/2)S$, $k_R<4S$, $2<e<3$ and $3<\pi<4$ give
\begin{equation}\label{eq:B-closed}
 B_R<300S^{-3}e^{-\pi S/2},\qquad
 L_R^*<\frac92+2\log S.
\end{equation}
For the first constant, $\pi^2/e^4<1$ and $e^{-\pi/2}<1$ reduce it
to $11520\cdot16/625=184320/625<300$. Since $\sqrt S>2$,
\[
 L_R^*+(8R+4)e^R<7(1+\log S)\sqrt S.
\]
The first term of \eqref{eq:eps} is therefore at most
$4200(1+\log S)S^{-5/2}e^{-\pi S/2}$. Finally
$S^{9/2}e^{-\pi S/2}$ decreases for $S\ge4$ and is less than $8$ there:
at $4$ it is $512e^{-2\pi}<512e^{-6}<8$. Thus
\eqref{eq:residual-closed} contributes less than another
$4800(1+\log S)S^{-5/2}e^{-\pi S/2}$. This proves
\eqref{eq:closed-eps} and completes the proof of
Theorem~\ref{thm:main}.

\section{Support matching and the remaining signed interior}

\begin{corollary}[Matched finite-prime localization]\label{cor:finiteN}
For an integer $N\ge81$, set
\begin{equation}\label{eq:matched-R}
 R=\tfrac12\log N-1\ge1,\qquad S=N/e^2.
\end{equation}
Then for $u\in\mathcal D_0$,
\begin{align}
 \left|\HH_N(u)-\HH_N(\chi_Ru)-t_\Gamma(\eta_Ru)\right|
       &\le(\varepsilon_R+r_N+\tau_N)N_X(u),\label{eq:finite-H}\\
 \left|q_N(u)-q_N(\chi_Ru)-t_\Gamma(\eta_Ru)\right|
       &\le(\varepsilon_R+r_N)N_X(u).\label{eq:finite-q}
\end{align}
\end{corollary}
\begin{proof}
The localized input $v=\chi_Ru$ is supported in
$[-\tfrac12\log N,\tfrac12\log N]$. Its correlation is zero for every
displacement $s\ge\log N$. At equality the possible overlap is a single
point and has measure zero. Hence, exactly,
\begin{equation}\label{eq:support-match}
 p_\infty(v)=p_N(v),\qquad T_N(v)=0,\qquad
 q_\infty(v)=q_N(v)=\HH_N(v).
\end{equation}
This does not discard a nonzero prime-power atom at $n=N$.
Use \eqref{eq:phase-form} on $u$ only and
\eqref{eq:support-match} on $v$, followed by Theorem~\ref{thm:main}.
For \eqref{eq:finite-q}, use
$|q_N(u)-q_\infty(u)|\le r_NN_X(u)$ instead. No second interior
truncation error is needed.
\end{proof}

All errors in the corollary tend to zero. The leading displayed bound has
the scale
\[
 (1+\log N)N^{-5/2}\exp\left(-\frac{\pi N}{2e^2}\right)
\]
up to fixed constants. For fixed $R$, any integer
$K\ge\lceil e^{2(R+1)}\rceil$ gives the exact identity
$q_\infty(\chi_Ru)=q_K(\chi_Ru)$. In both statements the full continuum
of gamma shifts is retained; even shifts larger than the support diameter
contribute through the difference norm in \eqref{eq:gamma}.

\begin{remark}[Centered inputs can acquire a negative pole defect]
If $\int_\R Cu=0$, then $\Pole(u)=0$, but multiplication by $\chi_R$
need not preserve that constraint. In \eqref{eq:exact-localization} the
pole defect is then exactly $-\Pole(\chi_Ru)\le0$. In particular
\eqref{eq:pole-error} gives
\[
 \Pole(\chi_Ru)\le
 \bigl[8e^RB_R+(1+S^{-1})\tau_S\bigr]N_X(u).
\]
The negative leakage has already been included in $\varepsilon_R$.
Dropping it on the ground that the original input is centered would be
incorrect.
\end{remark}

The signed difference in \eqref{eq:main} polarizes to an $X$-bounded
symmetric operator of norm at most $\varepsilon_R$. This statement concerns
the difference, not either unbounded gamma-containing form separately.
Likewise \eqref{eq:finite-H} is not a bound
$\HH_N(u)\ge-o(1)N_X(u)$: it retains the unresolved interior term
$\HH_N(\chi_Ru)$. No smallest positive interior margin is certified,
and no comparison of such a margin with the uniform error is asserted.

The arguments use only elementary weighted correlation estimates,
the theta relation proved below, and exact algebra. They do not use a
prime number theorem, a zero-location or zero-simplicity hypothesis,
a spectral projection, or a ground-state/radical identity. Their outcome
is an explicit vanishing localization defect with a signed
infinite-dimensional interior left intact. Neither the Riemann hypothesis
nor its negation follows.

\appendix
\section{Theta kernel and normalization}
\label{app:theta}

This appendix supplies the theta facts used in the localization argument,
including all normalization factors.  For $t>0$ and $x\in\mathbb R$, define
\begin{equation}\label{eq:app-theta-definitions}
 \begin{gathered}
 \vartheta(t)=\sum_{n\in\mathbb Z}e^{-\pi n^2t},\qquad
 G(x)=e^{x/2}\vartheta(e^{2x}),\qquad
 \kappa(x)=\frac{G''(x)-G(x)/4}{2},\\
 C(x)=\cosh(x/2),\qquad W(x)=\frac{\kappa(x)}{2C(x)}.
 \end{gathered}
\end{equation}
We prove that $\kappa$ and $W$ are smooth, even and strictly positive, that
\begin{equation}\label{eq:app-theta-envelope}
 0<W(x)\le 40\exp\!\left(4|x|-\pi e^{2|x|}\right),
\end{equation}
and that
\begin{equation}\label{eq:app-theta-mass}
 \int_{\mathbb R}2C(x)\kappa(x)\,\mathrm dx=1.
\end{equation}
In particular, \eqref{eq:app-theta-mass} is equivalently
$\int_{\mathbb R}4C(x)^2W(x)\,\mathrm dx=1$; it is not a claim that
$W$ itself has unit integral.

\paragraph{Gaussian periodization.}
For completeness, the needed theta transformation follows directly from
the Gaussian.  In the Fourier convention
\[
 F_t(\xi)=\int_{\mathbb R}e^{-\pi t y^2}e^{-2\pi i\xi y}\,\mathrm dy,
\]
differentiation under the integral and integration by parts give
$F_t'(\xi)=-(2\pi\xi/t)F_t(\xi)$.  Also $F_t(0)=t^{-1/2}$: indeed,
the square of $\int_{\mathbb R}e^{-\pi y^2}\,\mathrm dy$ is, by polar
coordinates,
$2\pi\int_0^\infty r e^{-\pi r^2}\,\mathrm dr=1$.
Solving the differential equation therefore gives
\begin{equation}\label{eq:app-gaussian-fourier}
 F_t(\xi)=t^{-1/2}e^{-\pi\xi^2/t}.
\end{equation}
All differentiations and integrations by parts here are justified by
integrable polynomial multiples of a Gaussian.

The periodization
$P_t(y)=\sum_{n\in\mathbb Z}e^{-\pi t(y+n)^2}$ is smooth and
$1$-periodic, with locally uniformly convergent differentiated series.
Its $k$th Fourier coefficient, by an absolutely convergent change of
variables, is $F_t(k)$.  Thus
\[
 P_t(y)=t^{-1/2}\sum_{k\in\mathbb Z}
                  e^{-\pi k^2/t}e^{2\pi i k y}.
\]
For an explicit uniqueness justification, both sides are continuous and
have the same Fourier coefficients; convolution of their difference with
the Fej\'er kernel
$K_N(y)=(N+1)^{-1}|\sum_{r=0}^Ne^{2\pi i r y}|^2$ is zero.
These convolutions converge uniformly to the difference: $K_N\ge0$,
$\int_0^1K_N=1$, and on every closed set away from the integers,
$K_N=O((N+1)^{-1})$.  Uniform continuity supplies the remaining
neighborhood of the integers.  Evaluating at $y=0$ proves
\begin{equation}\label{eq:app-theta-inversion}
 \vartheta(t)=t^{-1/2}\vartheta(1/t).
\end{equation}
Termwise differentiation of $\vartheta$ is valid on every compact
subinterval of $(0,\infty)$.  Differentiating
\eqref{eq:app-theta-inversion} at $t=1$ yields
\begin{equation}\label{eq:app-theta-derivative-one}
 \vartheta'(1)=-\frac14\vartheta(1).
\end{equation}

\paragraph{Smoothness, parity and positivity.}
Equation~\eqref{eq:app-theta-inversion} gives $G(-x)=G(x)$, so $G$,
$\kappa$, and $W$ are smooth and even.  Put $t=e^{2x}$.  Direct
differentiation of \eqref{eq:app-theta-definitions} gives
\begin{equation}\label{eq:app-theta-kappa-derivative}
 \kappa(x)=e^{x/2}\bigl(2t^2\vartheta''(t)+3t\vartheta'(t)\bigr).
\end{equation}
The constant term of $\vartheta$ is annihilated.  Pairing its positive and
negative indices gives, in particular for $x\ge0$,
\begin{equation}\label{eq:app-theta-kappa-series}
 \kappa(x)=\sum_{j\ge1}
 \left(4\pi^2j^4e^{9x/2}-6\pi j^2e^{5x/2}\right)
             e^{-\pi j^2e^{2x}}.
\end{equation}
For $x\ge0$, each coefficient is
$e^{x/2}(4u^2-6u)>0$, where $u=\pi j^2e^{2x}\ge\pi>3$.
Consequently $\kappa>0$ on that half-line, and evenness gives positivity
everywhere.  Since $C>0$, the same holds for $W$.

\paragraph{A complete uniform envelope.}
The elementary bounds $3<\pi<22/7$ can be obtained without numerical
input from
\[
 \pi=4\int_0^1\frac{\mathrm du}{1+u^2}
 >4\sum_{k=0}^7\frac{(-1)^k}{2k+1}
 =\frac{135904}{45045}>3,
 \qquad
 \frac{22}{7}-\pi
 =\int_0^1\frac{u^4(1-u)^4}{1+u^2}\,\mathrm du>0.
\]
The first inequality follows by the finite geometric expansion with
positive remainder, and the second identity by polynomial division.
Also $e^3>\sum_{k=0}^4 3^k/k!=131/8>16$, so
$e^{-3\pi}<e^{-9}<1/4096$.

For $x\ge0$, use $2C(x)\ge e^{x/2}$ and discard the negative term in
\eqref{eq:app-theta-kappa-series}.  Since
$j^2-1\ge3(j-1)$, this gives
\begin{align*}
 W(x)
 &\le4\pi^2e^{4x-\pi e^{2x}}
           \sum_{j\ge1}j^4e^{-\pi(j^2-1)e^{2x}}\\
 &\le4\left(\frac{22}{7}\right)^2e^{4x-\pi e^{2x}}
           \sum_{j\ge1}\frac{j^4}{4096^{j-1}}.
\end{align*}
This bounds the full theta tail, not merely its first summand.
For every integer $j\ge1$, $j^4\le16^{j-1}$, by induction using
$((j+1)/j)^4\le16$.  Hence the coefficient above is bounded by
\[
 4\left(\frac{22}{7}\right)^2
   \sum_{j\ge1}\frac{j^4}{4096^{j-1}}
 \le4\left(\frac{22}{7}\right)^2\frac{256}{255}
 =\frac{495616}{12495}<40,
 \qquad 40-\frac{495616}{12495}=\frac{4184}{12495}>0.
\]
Evenness proves \eqref{eq:app-theta-envelope} on the whole line.  In
particular, $2C\kappa=4C^2W$ is integrable.

\paragraph{Exact mass.}
For $t\ge1$, the differentiated theta series gives
$\vartheta(t)\to1$ and $t^a\vartheta'(t)\to0$ for every fixed real $a$.
For example, after factoring out $e^{-\pi t}$, the remaining sums are
bounded by the convergent series
$\sum_{j\ge1}j^m e^{-\pi(j^2-1)}$, with a suitable fixed $m$.
Using evenness, $\mathrm dx=\mathrm dt/(2t)$, and
\eqref{eq:app-theta-kappa-derivative}, one obtains
\begin{align*}
 \int_{\mathbb R}2C\kappa\,\mathrm dx
 &=\int_1^\infty
    \left[2(t^{3/2}+t)\vartheta''(t)
                +3(\sqrt t+1)\vartheta'(t)\right]\,\mathrm dt\\
 &=\left[2(t^{3/2}+t)\vartheta'(t)+\vartheta(t)\right]_1^\infty\\
 &=1-4\vartheta'(1)-\vartheta(1)=1,
\end{align*}
where the last equality is exactly
\eqref{eq:app-theta-derivative-one}.  The displayed theta tail bounds
justify the endpoint limits and the integrations.  This proves
\eqref{eq:app-theta-mass} directly, with no Mellin transform or
spectral assumption.

\section*{Data availability and manuscript status}
The companion manuscript source, theta appendix and source-free scalar
checks are provided at
\url{https://github.com/serabing-hash/theta-weighted-certificate/tree/main/research/2026-10-27-signed-cubic-localization}.
The scalar checks verify recorded finite arithmetic; they do not certify
the analytic proof. Artificial intelligence assistance was used for
drafting, mathematical exposition checks and document preparation.
This manuscript is a preprint and makes no claim of human peer review.

\end{document}
